% 02_langage_et_backend.tex
\subsection{Évolution de l'expressivité de PlanningSpec}

Lors de notre évaluation initiale du métamodèle PlanningSpec par rapport à l'implémentation MiniZinc, nous avons identifié deux écarts sémantiques majeurs et trois faiblesses d'ordonnancement par rapport au modèle général du RCPSP. 

\subsubsection{Corrections Sémantiques (BugFixes)}
\begin{itemize}
    \item \textbf{La préférence $MaxPerScope$} : L'attribut \texttt{resourceType} était absent de la grammaire Langium. En conséquence, la limitation portait sur le nombre d'activités globales simultanées, et non sur le \textit{Workload Balance} des ressources. Nous avons réinséré l'attribut et corrigé le générateur en \texttt{sum(r in RES) (bool2int(...))}.
    \item \textbf{La pénalité $Day-Off$} (\texttt{AvoidParticipationOnDate}) : La pénalité n'était pas appliquée de façon unitaire par jour travaillé, mais multipliée par le nombre de créneaux assurés par la ressource ce jour-là. Le générateur a été corrigé via une borne booléenne \texttt{bool2int(sum(...) > 0) * weight}.
\end{itemize}

\subsubsection{Nouvelles Règles Structurelles}
Afin d'étendre la puissance de la planification aux cas d'usage académiques avancés (ex: Jurys, Délibérations, Soutenances), nous avons étendu la grammaire Langium :

\begin{enumerate}
    \item \textbf{MandatoryRoles (Hard)} : 
    Pour s'assurer qu'activités complexes (soutenances) n'abandonnent pas leurs rôles, cette contrainte oblige le solveur à trouver (au moins) une ressource pour chaque rôle énuméré : \texttt{sum(bool2int(role\_assignment)) >= 1}.
    
    \item \textbf{TemporalPrecedence (Hard)} :
    Emprunt direct au PERT / RCPSP. Oblige l'activité $A$ à terminer avant le début de $B$ : \texttt{start\_time[A] + duration[A] <= start\_time[B]}.
    
    \item \textbf{TimeWindow (Hard)} :
    Astraint une instance d'activité à s'exécuter dans des bornes temporelles fixes : \texttt{minSlot $\leq$ start\_time} et \texttt{end\_time $\leq$ maxSlot}.
    
    \item \textbf{PreferredResource (Soft)} :
    Alors que \texttt{FixedAssignment} force une ressource ("Il DOIT être le président du jury"), \texttt{PreferredResource} exprime un souhait pénalisant s'il n'est pas honoré : \texttt{penalty = weight * not(role\_assignment)}.
\end{enumerate}

Le métamodèle PlanningSpec est ainsi rendu strictement conforme à l'expression de tout problème d'ordonnancement RCPSP, en étendant ses capacités aux plannings académiques multicritères.
